\section{Colorings of knot diagrams}\label{sec:knotcolor}
The figure-eight knot pictured at the beginning of \emph{Topology} has
a four-crossing diagram. Can one turn this diagram into a circle by
moving its strands? A count with five colors answers the question.
A separate four-letter computation identifies the knot with its mirror.

\subsection{Local rules and Reidemeister moves}
Use a finite diagram with crossing records. Each crossing has four
half-edges in cyclic order, an opposite pair marked over, and the other
pair marked under. Edges pair the half-edges between crossings. Store
each crossing-free circle separately. For planar diagrams one may give a
planar drawing or a list of elementary tangle slices as well. An
$n$-coloring, $n\ge1$, assigns each edge (and each free circle) a residue
modulo $n$. At a crossing the two over colors agree, say with $y$, and the
two under colors $x,x'$ satisfy $x+x'=2y$. Thus the under color changes
by the operation $x*y=2y-x$. Following a strand across successive
overcrossings gives the equivalent familiar coloring of arcs between
undercrossings. Write $C_n(D)$ for the number of colorings.

\begin{theorem}[Fox coloring invariance]\label{idx-knots20}
A Reidemeister move induces a bijection on $n$-colorings fixing the
untouched part of the diagram. In particular $C_n$, and the existence of
a nonconstant coloring, are invariant under a finite sequence of
diagram moves, for every positive modulus.
\end{theorem}
\begin{proof}
In a type I kink the over color is also one of the under colors,
since the small loop directly joins the corresponding half-edges. The
crossing equation forces the other under color to agree. The local
extension is unique and has equal colors at both ends.
In type II one strand passes over twice with the same color $y$. A
color $x$ on the other strand changes first to $x*y$ and then back to
$x$, since $(x*y)*y=x$. This determines all internal colors uniquely.

For type III name the strands by their relative height. Read the top
color as $z$, the middle color on the side before it passes under the
top as $y$, and the bottom input as $x$. On one side of the move the
bottom passes under the middle and then the top, leaving with
$(x*y)*z$. On the other side it passes under the top and then under
the transformed middle, leaving with $(x*z)*(y*z)$. These agree since
both equal $x-2y+2z$. The top stays $z$ and the middle exits with $y*z$
on both sides. These three parameters determine all local colors, so
the two local assignments agree at every exterior port. Directions here
are choices for tracing the three strands. Reading a strand backwards
uses the same crossing relation because it is involutive. The argument
therefore also covers reversed and rotated local pictures. Gluing the
unchanged exterior identifies the full colorings even if some exterior
ports are joined to one another. Constant colorings extend constantly
in both directions. Relabeling the edges simply permutes color
coordinates. Composing the local bijections proves the assertion.
\end{proof}

\subsection{Two-strand links and the figure-eight}
A braid word is read from top to bottom, and its closure joins each bottom
position to the top position of the same index. For a positive crossing
with the left strand passing over the right, the pair of colors transforms
by
\[
 B(x,y)=(2x-y,x),\qquad B^{-1}(x,y)=(y,2y-x).
\]
All untouched positions keep their colors. Once the top colors are
given, all internal colors are forced, so the colorings of the closure
are exactly the fixed vectors of the braid word modulo $n$.

\begin{theorem}[Two-strand closure formula]\label{idx-knots65}
For the closure $D_m$ of the two-strand braid with signed exponent
$m\in\mathbb Z$,
\[
                  C_n(D_m)=n\gcd(n,|m|).
\]
For odd $m$ this is a knot, and for even $m$ it is a two-component link.
The positive odd exponents give pairwise distinct diagram classes.
\end{theorem}
\begin{proof}
Put $a=x-y$. Both crossings preserve this difference. The positive
crossing adds $a$ to both entries and its inverse subtracts $a$.
Iteration gives $B^m(x,y)=(x+ma,y+ma)$. The closure condition is thus
$ma=0$ modulo $n$. Set $h=\gcd(n,|m|)$ and $n'=n/h$.
The admissible residues are exactly $0,n',\ldots,(h-1)n'$.
Indeed, cancelling $h$ gives $n'\mid (m/h)a$, and B\'ezout's identity
for the coprime integers $n',m/h$ gives $n'\mid a$. The converse
follows by multiplication. The identity is obtained by the ordinary
Euclidean algorithm with decreasing remainders, carrying the two
coefficients in each division. For $m=0$ every residue works. There are
$n$ independent choices of $y$, giving the formula. Every crossing
transposes the two strand positions, so the cycles of the endpoint
permutation give the asserted component counts. If $1\le m<k$ are odd,
use $n=k$: $C_k(D_k)=k^2$ while $C_k(D_m)\le km<k^2$.
\end{proof}

The trefoil is the closure $D_3$ and has nine three-colorings, six
nonconstant. At modulus three the crossing equation says that the
three arc colors are all equal or all different, as checking $x=y$ or
$x\ne y$ shows. The crossing-free unknot has exactly three
three-colorings, one per constant color. Thus no finite Reidemeister
sequence takes one of these two diagrams to the other.

\begin{proposition}[The figure-eight diagram]\label{idx-knots98}
Write $a=\sigma_1,b=\sigma_2$ for the adjacent crossings of three
strands. The figure-eight diagram is the closure of $w=(ab^{-1})^2$.
It has $C_n(w)=n\gcd(n,5)$. Its diagram class is distinct from the
unknot and the trefoil. It is equivalent to its mirror and has
unknotting number one under diagram moves and crossing changes.
\end{proposition}
\begin{proof}
Apply $B$ on slots $1,2$, then $B^{-1}$ on slots $2,3$, twice.
For input $(x,y,z)$, put $u=x-y$, $v=y-z$. Subtracting the input from
the output gives $(3u+v,-u-2v,-4u-3v)$. Vanishing is equivalent to
$v=-3u$ and $5u=0$ modulo $n$. Given $z,u$, these equations reconstruct
all three input colors. Counting $u$ as in the previous proof gives the
formula. At $n=5$ the count is $25$, whereas the trefoil and unknot
both have $5$. At $n=3$ the figure-eight has only its constant colorings.

Mirroring across the projection plane changes all crossing signs, giving
$(a^{-1}b)^2$. The braid relation $aba=bab$ is a type III move. Thus
$\Delta=aba=bab$ satisfies $\Delta a=b\Delta$ and $\Delta b=a\Delta$ by
concatenation. Conjugation by $\Delta$ interchanges $a,b$, including
their inverses. Conjugation preserves a closed braid diagram class:
move the cut around the annular closure, which cyclically permutes the
word, then cancel inverse pairs by type II. Hence the mirror has the
class of $(b^{-1}a)^2$, which cyclically shifts to $w$.

Change the first $a$ in $w$ to $a^{-1}$. Cyclically shifting its last
letter to the front gives $b^{-1}a^{-1}b^{-1}a$. The inverse braid
relation changes the first three letters to $a^{-1}b^{-1}a^{-1}$, so
cancellation leaves $a^{-1}b^{-1}$. Its rightmost strand now takes
part in only its last crossing. On closure the strand and its outer
closing arc form an empty kink, removable by a type I move after
straightening that arc. The two-strand closure left over is $D_{-1}$,
again a single removable kink. One crossing change therefore suffices.
Zero cannot suffice by the count at modulus five.
\end{proof}

\subsection{Joining knots and the arithmetic statement}
For connected sums use diagrams cut into $(1,1)$-tangles, each drawn
with one input at the top and one output at the bottom. Closing the cut
by an exterior arc gives its original diagram. Stacking two tangles and
closing the result gives their marked connected sum. Represent a tangle
by successive horizontal slices containing one crossing, one birth of a
pair at a local minimum when read in the chosen direction, or one death
of a pair. These are combinatorial instructions about adjacent positions.

\begin{proposition}[Connected-sum count]\label{idx-knots143}
For two such presented knots,
\[
 n C_n(D_1\mathbin{\#}D_2)=C_n(D_1)\,C_n(D_2).
\]
In particular the connected sum of $s\ge1$ trefoils has $3^{s+1}$
three-colorings. Distinct positive values of $s$ give inequivalent
diagrams.
\end{proposition}
\begin{proof}
On a horizontal slice with colors $c_1,\ldots,c_k$ the alternating sum
$c_1-c_2+c_3-\cdots$ is conserved. A crossing preserves the difference
of the two entries it acts on. Inserting or removing an adjacent equal
pair contributes zero and changes no subsequent sign. At the two ends
there is only one entry, so every coloring has equal endpoint colors.
It extends uniquely across the closing arc. Translating all residues by
a constant preserves the equations and bijects the colorings at any two
specified endpoint colors. Thus for $D_i$ let $q_i$ be the number with
endpoint color zero. Its total is $nq_i$. For the stack, choose the
common endpoint color in $n$ ways and the two restricted colorings in
$q_1q_2$ ways. This gives the identity and, by induction from
$C_3(D_3)=9$, the claimed powers of three.
\end{proof}

For Lean, encode a crossing record by four edge indices and a bit for
the over pair. Define \texttt{colorOK} by one modular test and one
equality per crossing. With $e$ edge variables including free circles,
enumerate the $n^e$ vectors by radix digits (handle $n=1$ by the unique
zero vector). For each external port table the type I--III proofs give
explicit inverse functions between extension lists. Define
\texttt{braidStep} on indexed color lists, and prove by induction on
the word the fixed-vector test for the closure. Implement the signed
two-strand formula as two inductions on $|m|$. For the connected-sum
proof the induction predicate is equality of the current and initial
alternating sums, evaluated modulo $n$. Thus iPRA proves the invariant
for every coded move list, and unequal counts rule out any such list.
Classically Reidemeister's theorem connects this diagram equivalence to
ambient isotopy of tame knots. The iPRA claims above work directly with
diagrams and the displayed moves. A translation of arbitrary ambient
isotopies into move certificates would be a separate geometric theorem.
